% PCT-SOURCE: pdf=016 print=004 kind=chapter-title id=ch1-relativistic-transformation-laws
\section{Relativistic Transformation Laws}\label{ch:relativistic-transformation-laws}

% PCT-SOURCE: pdf=016 print=004 kind=epigraph id=epigraph-wigner-care
\begin{flushright}
\begin{minipage}{0.55\textwidth}
\small
``You point out that care is needed in the analysis of the representations
of the Lorentz group; I promise you that I will be careful.''

\medskip
\textsc{E. Wigner}
\end{minipage}
\end{flushright}

% PCT-SOURCE: pdf=016 print=004 kind=prose id=p4-heisenberg-picture
Throughout this book, states will be described in the Heisenberg picture of
quantum mechanics. The Schr\"odinger picture is much less convenient for the
description of a relativistic theory because it treats the time coordinate on
a very different footing from the space coordinates; as will be proved in
Chapter \ref{ch:general-theorems-relativistic-qft}, the other commonly used picture, the interaction picture, in
general does not exist. In the Heisenberg picture, to each state of the system
under consideration there corresponds a unit vector, say \(\ket{\Phi}\), in a
Hilbert space, \(\mathcal H\). The vector does not change with time, whereas
the observables, represented by hermitian linear operators acting on
\(\mathcal H\), in general do. The scalar product of two vectors \(\ket{\Phi}\)
and \(\ket{\Psi}\) in \(\mathcal H\) is denoted by
\(\braket{\Phi}{\Psi}\), called the \emph{transition amplitude} of the
corresponding states.

% PCT-SOURCE: pdf=016 print=004 kind=prose id=p4-unit-ray
Two vectors that differ only by multiplication by a complex number of modulus
one describe the same state, because the results of all experiments on a state
described by \(\ket{\Psi}\) may be expressed in terms of the quantities

% PCT-SOURCE: pdf=016 print=004 kind=display id=eq-p4-transition-probability
\begin{equation*}
  \left|\braket{\Phi}{\Psi}\right|^2
\end{equation*}

% PCT-SOURCE: pdf=016 print=004 kind=prose id=p4-unit-ray-cont
which gives the probability of finding \(\ket{\Phi}\) if \(\ket{\Psi}\) is
what you have. The set \(\{e^{i\alpha}\ket{\Phi}\}\) of vectors, where
\(\alpha\) varies over all real numbers, and the norm of \(\ket{\Phi}\)
(written \(\lVert\ket{\Phi}\rVert\)
and defined as \([\braket{\Phi}{\Phi}]^{1/2}\)) is unity, is called a
\emph{unit ray}. For brevity, we shall speak of the state \(\ket{\Phi}\). The
condition \(\lVert\ket{\Phi}\rVert=1\) is obviously equivalent to the
convention of normalizing the probability to unity. The preceding remarks
can be summarized: \emph{states}\footnote{%
% PCT-SOURCE: pdf=016 print=004 kind=footnote id=fn-p4-state
By ``state'' we shall always mean pure state. A ``mixed state'' can always be
formed from several states by a \emph{classical superposition}, each entering
with a certain known probability which describes our ignorance of the system.}
of a physical system are represented by unit rays.
